%%%%%%%%%%%%%%%%%%%%%%%%%%%%%%%%%
%
%       EXERCÍCIO
%
%%%%%%%%%%%%%%%%%%%%%%%%%%%%%%%%%

\ifdefstring{\atividade}{prova}{%
    \renewcommand{\valorquestao}{\ValorQ\ pontos}
}%

\begin{exercicioBanco}[\valorquestao]
    A função conjunta de probabilidade entre as variáveis \(X\) e \(Y\) é apresentada abaixo.
    %
    \[
    \renewcommand{\arraystretch}{1.5}
    \begin{array}{|c|c|c|c|c|c|}
        \hline
        X\backslash Y & -1 & 0 & 2 & 4 & P(X = x) \\
        \hline
        -2 &  & \frac{3}{64} & \frac{1}{32} &  & \frac{5}{16} \\
        \hline
        -1 & \frac{1}{16} & \frac{1}{16} & 0 &  &  \\
        \hline
        1 & \frac{1}{64} & \frac{11}{64} &  & \frac{1}{64} & \frac{5}{16} \\
        \hline
        2 & \frac{5}{64} &  & \frac{3}{64} & \frac{1}{32} &  \\
        \hline
        P(Y = y) &  & \frac{5}{16} &  & \frac{1}{4} & 1 \\
        \hline
    \end{array}
    \]
    %
    \begin{enumerate}[(a)]
        \item Complete a tabela.
        \item Obtenha as marginas de \(X\) e \(Y\).
        \item Calcule a função de probabilidade da variável \(XY\).
    \end{enumerate}
\end{exercicioBanco}

